\section{Quotient geometry and arithmetic witnesses}
\label{sec:quotient-witnesses}

\subsection{The exact nullspace}

\begin{theorem}[Completion quotient seminorm]
\label{thm:quotient-seminorm}
On \(\Ran R\), the optimal completion cost
\[
 \gamma(k)=\min_{RF=k}\|AF\|_1
\]
is a seminorm and
\begin{equation}
 \boxed{\ker\gamma=R(\ker A).}
 \label{eq:gauge-nullspace}
\end{equation}
Consequently it induces a norm on
\[
 \Ran R\big/R(\ker A).
\]
\end{theorem}

\begin{proof}
Positive homogeneity follows by scaling a completion.  If
\(RF_i=k_i\), then \(R(F_1+F_2)=k_1+k_2\) and
\[
 \|A(F_1+F_2)\|_1
 \le \|AF_1\|_1+\|AF_2\|_1.
\]
Taking minima gives the triangle inequality.

If \(k=RF\) for some \(F\in\ker A\), then \(\gamma(k)=0\).
Conversely, if \(\gamma(k)=0\), primal attainment in
\cref{thm:strong-completion-duality} supplies \(F\) with
\(RF=k\) and \(\|AF\|_1=0\), so \(F\in\ker A\).
\end{proof}

\begin{corollary}[Stability]
\label{cor:gauge-stability}
For \(k,k'\in\Ran R\),
\[
 |\gamma(k)-\gamma(k')|
 \le \gamma(k-k').
\]
Thus the gauge is Lipschitz with constant one in its own quotient norm.
\end{corollary}

\begin{proof}
Apply the triangle inequality twice.
\end{proof}

\subsection{The dual unit ball}

Define
\[
 \cD_{A,R}
 =
 \left\{
 \lambda\in\cK:
 \exists z,\ A^*z=R^*\lambda,\ \|z\|_\infty\le1
 \right\}.
\]
When \(R\) is surjective this set is compact, convex, balanced, and
\begin{equation}
 \gamma(k)
 =
 \max_{\lambda\in\cD_{A,R}}
 [\lambda,k]_{\mathbb R}.
 \label{eq:gauge-support-function}
\end{equation}
It is the unit ball of the dual of the quotient norm in
\cref{thm:quotient-seminorm}.  A lower obstruction is therefore not a
heuristic roughness statistic; it is a separating functional in the exact
dual unit ball.

\begin{proposition}[Zero-cost criterion by annihilators]
\label{prop:zero-cost-annihilator}
For \(k\in\Ran R\), the following are equivalent:
\begin{enumerate}[label=\textup{(\roman*)},leftmargin=2.6em]
 \item \(\gamma(k)=0\);
 \item \(k\in R(\ker A)\);
 \item \([\lambda,k]_{\mathbb R}=0\) for every
       \(\lambda\) satisfying
       \(R^*\lambda\in\Ran A^*\).
\end{enumerate}
\end{proposition}

\begin{proof}
The first two statements are
\cref{thm:quotient-seminorm}.  In finite dimensions,
\[
 \bigl(R(\ker A)\bigr)^\perp
 =
 \{\lambda:R^*\lambda\in(\ker A)^\perp\}
 =
 \{\lambda:R^*\lambda\in\Ran A^*\}.
\]
The third statement says exactly that \(k\) belongs to the double
orthogonal complement of \(R(\ker A)\).
\end{proof}

\subsection{The physical M\"obius direction}

For the primitive lattice \(P\), define
\[
 \sigma(u,v)=\mu(u)\mu(v)
 \qquad((u,v)\in P).
\]
The physical functional is
\[
 \cS_{\rm prim}(k)=\langle\sigma,k\rangle.
\]

\begin{theorem}[Arithmetic dual witness]
\label{thm:arithmetic-dual-witness}
For every \(D\ge1\), there is
\(z^\sigma_D\in\cY_D\) such that
\begin{equation}
 A_D^*z^\sigma_D=R^*\sigma,
 \qquad
 \|z^\sigma_D\|_\infty\le1.
 \label{eq:arithmetic-dual-feasible}
\end{equation}
Consequently,
\begin{equation}
 \boxed{
 |\cS_{\rm prim}(k)|
 \le \gamma_D(k).}
 \label{eq:physical-below-gauge}
\end{equation}
\end{theorem}

\begin{proof}
The exact diagonal divisor lift followed by Abel summation
\citep{WangTPC71} gives, for every completion \(F\),
\begin{align*}
 \langle\sigma,RF\rangle
 ={}&
 \sum_{d\le D}\mu(d)
 \sum_{x,y}M_d(x)M_d(y)
 \Delta_{12}F^{[d]}(x,y)\\
 &+
 \sum_{d>D}\mu(d)
 \sum_{x,y}\mu(dx)\mu(dy)F(dx,dy).
\end{align*}
Explicitly, with \(\operatorname{sgn}(0)=0\), take
\begin{align*}
 (z^\sigma_D)^{\rm low}_{d,x,y}
 &=
 \operatorname{sgn}\!\left(
 \mu(d)M_d(x)M_d(y)
 \right),\\
 (z^\sigma_D)^{\rm tail}_{d,x,y}
 &=
 \operatorname{sgn}\!\left(
 \mu(d)\mu(dx)\mu(dy)
 \right).
\end{align*}
Thus every nonzero coordinate is the corresponding real coefficient
divided by its absolute value.  This defines \(z^\sigma_D\) with coordinate
modulus at most one and
\[
 \langle z^\sigma_D,A_DF\rangle
 =
 \langle\sigma,RF\rangle
 \qquad(F\in\cX).
\]
Hence \(A_D^*z^\sigma_D=R^*\sigma\).

For any unit complex number \(\omega\), the pair
\((\omega\sigma,\omega z^\sigma_D)\) is also dual feasible.  Choose
\(\omega\) so that
\(\Real\langle\omega\sigma,k\rangle
=|\langle\sigma,k\rangle|\), and apply
\cref{cor:tpc73-gauge-dual}.
\end{proof}

\begin{remark}[The witness need not be optimal]
The arithmetic sign witness recovers the original TPC--73 upper
certificate for the physical functional.  The dual optimizer may give a
strictly larger value than
\(|\cS_{\rm prim}(k)|\).  Such a gap obstructs the chosen
absolute-value certificate architecture; it does not prove that the
physical signed functional is large.
\end{remark}

\begin{corollary}[Single-atom rigidity]
\label{cor:single-atom-rigidity}
If \(k=ce_{u_0,v_0}\) with
\((u_0,v_0)=1\) and \(\mu(u_0)\mu(v_0)\ne0\), then
\[
 \gamma_D(k)\ge|c|.
\]
No nonprimitive analytic fill can erase this nonzero physical atom.
\end{corollary}

\begin{proof}
Use \cref{eq:physical-below-gauge}.
\end{proof}
